\section{Actual first-order components have linear positive-dimensional degree}
\label{app:actual-first-order}

The degree of a Taylor-reconstruction cover can exceed the degree of
the actual polynomial-solution locus. The following bound separates
its curve and surface components. We then exhibit quadratically many
actual isolated regular solutions, even at derivative degree one,
and prove that this example still has only constantly many nearby
labels at a fixed positive agreement gap. The full first-order MCA
conjecture remains unresolved.

\begin{theorem}[Positive-dimensional actual components]
\label{thm:actual-first-order-components}
Let $E$ be algebraically closed of characteristic zero or $p>D$, with
$D\ge1$. Let $Q(X,z,u,v)\ne0$ have total degree at most $B\ge1$ in
$(u,v)$ and degree at most $H$ in $z$, with no restriction on its
$X$-degree. In affine $(z,P)$ coefficient space, let $\mathcal V$ be
the Zariski closure of the actual regular solutions
\[
 \deg P\le D,\qquad Q(X,z,P,P')=0,\qquad
 Q_v(X,z,P,P')\not\equiv0.
\]
Put
\[
 \tau=\max\{0,2D-3\},\quad b=1+\tau(B-1),\quad
 T=b+\tau H,\quad K=H+BT,\quad q=B+H.
\]
Every component of $\mathcal V$ has dimension at most two. If $J_i$
is the sum of the projective degrees of its $i$-dimensional irreducible
components, then
\[
 J_1\le qK,\qquad J_2\le qT.
\]
In particular $J_1+J_2=O_{B,H}(D)$. The reduced generic challenge
fiber of the surface components has total degree at most $\deg_v Q$.
The bounds concern actual components defined over $E$ before choosing
a generic Taylor center; they do not apply to arbitrary Taylor images.
\end{theorem}
\begin{proof}
If the regular locus is empty, the claims are immediate. Otherwise
adjoin a Taylor center $t$ transcendental over $E$. Put
$S=Q_v$, $A_0=Q_X+vQ_u$, and
\[
 \mathcal D=\partial_X+v\partial_u-(A_0/S)\partial_v,
\]
holding $z$ constant. On $S\ne0$, implicit differentiation reconstructs
every derivative through order $D$. For $j\ge2$, write this derivative
as $N_j/S^{2j-3}$, with $N_2=-A_0$. If $e=2j-3$, the recurrence is
\[
\begin{split}
 N_{j+1}={}&S^2(\partial_XN_j+v\partial_uN_j)
       -A_0S\partial_vN_j\\
 &-eN_j\bigl(S(\partial_XS+v\partial_uS)-A_0\partial_vS\bigr).
\end{split}
\]
Induction gives jet degree at most $1+(2j-3)(B-1)$ and challenge
degree at most $(2j-3)H$. The common-denominator Taylor polynomial
through degree $D$ is $F(X)/S(t,z,u,v)^\tau$, where every coefficient
of $F$ has jet degree at most $b$ and challenge degree at most $\tau H$.
Division by the Taylor factorials is valid since $p>D$.

\paragraph{The actual locus in initial data.}
Substitute this polynomial and its derivative into $Q$ and clear
$S^{\tau B}$. In a term of jet degree $e\le B$, the numerator has
jet degree at most
\[
 e b+\tau(B-e)(B-1)\le B b
\]
and challenge degree at most $H+\tau BH$. Thus every $X$-coefficient
of the residual has total $(z,u,v)$ degree at most $K$.
Include the initial equation $Q(t,z,u,v)=0$, of total degree at most $q$.
On $S(t,z,u,v)\ne0$, these equations describe exactly the initial
data of actual degree-$\le D$ solutions. Reconstruction and the
initial-jet map are inverse regular maps on these open loci.
Consequently all actual regular components have dimension at most two.
Their closures acquire no irreducible component supported entirely
on the boundary of this open locus.

\paragraph{Curve components.}
Choose $\lambda,c$ independently generic after adjoining $t$, and
cut the actual locus by
\[
 P(t)+\lambda z=c.
\]
Each projective curve component has finitely many points at infinity,
defined over $E$. At any such point the coefficient direction and
the $z$ direction are not all zero, so
$\sum_i a_i t^i+\lambda z\ne0$. The chosen hyperplane avoids these
points. Its generic intersection also avoids the finite separant
boundary and the finite intersections of each curve with other
actual components, including surfaces.

The function $P(t)+\lambda z$ on each curve is separable. The
coefficient functions and $z$ generate the curve's function field,
so their differentials cannot all vanish; independence of $t,\lambda$
prevents cancellation. Its generic fiber is therefore reduced.
The section selects exactly $J_1$ distinct affine points in the
regular reconstruction locus, away from the surface components.

In initial data the section substitutes $u=c-\lambda z$ into
$Q(t,z,u,v)$, a nonzero plane equation of degree at most $q$.
Each residual restricts to degree at most $K$. A plane component
on which all residuals vanish, unless contained in $S=0$, maps to
a positive-dimensional part of this generic hyperplane section of
the actual locus, and hence to a section of a surface component.
None of the selected $J_1$ points lies on it. Nor do they lie on
components contained in $S=0$. On every remaining plane component
at least one residual is nonzero. A generic linear combination of
the residuals vanishes identically on none of them. B\'ezout's bound
for its intersections with their union gives $J_1\le qK$.

\paragraph{Surface components.}
Use the same generic section. It contains no curve at infinity of
an actual surface: a generic form $P(t)+\lambda z$ is nonzero on
each such curve. Thus the affine section closures retain the full
surface degrees. On the regular initial-data locus the projection
to $(z,u)$ is generically separable, because $Q_v\ne0$. A generic
line $u+\lambda z=c$ therefore gives generically reduced sections.

Their initial-data curves are distinct components of the plane
equation $Q(t,z,c-\lambda z,v)=0$ and have summed degree at most $q$.
The projective reconstruction coordinates can be written as
\[
 S^\tau,\quad zS^\tau,\quad\text{the coefficient numerators of }F.
\]
After the substitution, all have degree at most $T$: the first two
have degrees at most $T-1,T$, and the remaining numerators have
degree at most $b+\tau H=T$. Homogenizing to degree $T$, the
degree of the rational image of each source curve is at most $T$
times its degree, as seen by pulling back a generic hyperplane.
Summing gives $J_2\le qT$.

Finally fix the generic challenge and work over its algebraic
closure. At an independent generic Taylor center, regular
reconstruction identifies every positive-dimensional actual
solution component with a component of the plane equation in
$(u,v)$. Projection to $u$ has total degree at most $\deg_v Q$.
The hyperplane $P(t)=u_0$ avoids all curve points at infinity by
the same transcendence argument, and its generic fiber is reduced.
It therefore counts the full coefficient-space degree of these
curves, proving the generic-fiber assertion.
\end{proof}

\paragraph{Linear degree cannot be replaced by a constant.}
Let $D\ge2$, in characteristic zero or $p>D+1$, and put
\[
 R(X,z)=X^{D+1}+zX-1,\qquad
 Q=RY_1-R_XY_0+Y_0^2.
\]
Here $B=2$, $H=1$. Every nonzero polynomial solution satisfies
$(R/P)'=1$. The rational function $R/P-X$ has degree at most
$D+1$ and derivative zero, so it is constant. Indeed, in positive
characteristic its degree is less than $p$, whereas a nonconstant
rational $p$th power has degree at least $p$.
Hence all nonzero solutions are
\[
 z=a^{-1}-a^D,\qquad
 P_a(X)=\sum_{i=0}^{D-1}a^iX^{D-i}+a^{-1},\qquad a\ne0.
\]
Conversely these are solutions, since $R(a,z)=0$ and $P_a=R/(X-a)$.
The coefficient of $X^{D-1}$ recovers $a$. After multiplying the
projective coordinates by $a$, they span all monomials
$1,a,\ldots,a^{D+1}$, so the curve has degree $D+1$.
The additional zero-solution line gives $J_1=D+2$.

Linear surface growth also occurs: $(X-z)P'-DP=0$ has the actual
family $P=c(X-z)^D$. Two generic affine hyperplanes restrict to
$A(z)+cB(z)=0$ and $C(z)+cE(z)=0$, with $A,C$ of degree at most
one and $B,E$ of degree at most $D$. Their elimination polynomial
$AE-CB$ generically has degree $D+1$ and no extraneous roots.
Thus this surface has degree $D+1$, although its generic challenge
fiber is a line.

\begin{corollary}[Nonaffine actual curves contribute linearly]
\label{cor:nonaffine-actual-curves}
Fix a received line $f+zg$ on $n$ distinct evaluation points and an
agreement threshold $D<A\le n$. Among the actual curve components in
Theorem~\ref{thm:actual-first-order-components}, omit those that are
affine codeword graphs $P=F+zG$, with $\deg F,\deg G\le D$.
The remaining components contain at most
\[
 \frac{n-D}{A-D}\,qK
\]
polynomial--challenge pairs with at least $A$ agreements. In particular,
their exceptional-label contribution is $O_{B,H,\eta}(n)$ when
$A-D\ge\eta n$ and $D\le n$.
\end{corollary}
\begin{proof}
On an irreducible curve $C$, call a coordinate persistent if its
agreement hyperplane contains $C$. If more than $D$ coordinates were
persistent, interpolation at $D+1$ of them would force
$P=F+zG$ identically on $C$. A vertical curve would then be a point;
otherwise $C$ would be that affine codeword graph. Thus every retained
curve has at most $h\le D$ persistent coordinates.

Each of its other $n-h$ agreement hyperplanes meets $C$ in at most
$\deg C$ distinct affine points. Every selected nearby pair lies on
at least $A-h$ of these hyperplanes. Its number is therefore at most
\[
 (\deg C)\frac{n-h}{A-h}
 \le(\deg C)\frac{n-D}{A-D}.
\]
Sum over the retained curves and use $J_1\le qK$.
\end{proof}

\paragraph{Scope and verification of the component bounds.}
The theorem gives no linear bound on the number of isolated joint
solutions; the next construction shows that such a bound is false.
Nor does it bound by a constant the number of persistent
affine codeword graphs that can each contribute accidental agreements.
These obstacles prevent an inference of general $O(n/|F|)$ first-order
MCA from the component-degree bounds alone.
The checks in \path{research/first_order_actual_components/} verify
the reconstruction recurrence on five equations, $576$ finite Taylor
fixtures, and $40$ sharp-order linear-equation curves. A second checker
exhausts $33{,}849$ polynomial--challenge pairs for the nonlinear
curve family, verifies $39$ split reduced hyperplane sections, and
detects extra solutions outside that example's stricter characteristic
hypothesis $p>D+1$.
These checks supplement the geometric proof; they do not certify
the generic-section argument by themselves.

\subsection{Quadratically many isolated regular first-order solutions}
\label{sec:quadratic-isolated-first-order}

\begin{theorem}[Quadratic isolated count at derivative degree one]
\label{thm:quadratic-isolated-first-order}
Let $D\ge1$, in characteristic zero or $p>D+2$, and let $R$ be a
monic squarefree polynomial of degree $N=D+1$ with nonzero roots.
The first-order equation
\begin{equation}
 (z-X^2)(RP'-R'P+P^2)+2XRP-2R^2=0
 \label{eq:quadratic-isolated-riccati}
\end{equation}
has exactly $\binom{D+2}{2}$ polynomial--challenge solutions of
degree at most $D$, all regular. They are indexed by unordered
multisets $\{a,b\}$ of roots of $R$, including $a=b$, and are
\[
 z=ab,\qquad P_{a,b}=\frac{R}{X-a}+\frac{R}{X-b}.
\]
Thus its reduced actual joint solution locus consists of exactly
that many isolated points. Its jet degree is two, derivative degree
one, and challenge degree one. For every $D$, these points can all
have distinct labels over a prime field of size $O(D^3)$.
\end{theorem}
\begin{proof}
Put $U=(X^2+z)P-2XR$ and $V=R-XP$. Direct differentiation gives
\[
 U'V-UV'=(z-X^2)(RP'-R'P+P^2)+2XRP-2R^2.
\]
Since $R(0)\ne0$, $V$ is not identically zero. The rational function
$U/V$ has degree at most $D+2$ and derivative zero, so it is
constant. In positive characteristic this uses the strict bound
$D+2<p$: a nonconstant rational $p$th power has degree at least $p$.
Writing the constant as $-u$ gives
\[
 (X^2-uX+z)P=(2X-u)R.
\]
For $H=X^2-uX+z$, this says $P=RH'/H$. Every root of $H$ must
belong to $R$: a root of multiplicity one or two gives a nonzero
simple-pole residue in $H'/H$, since the characteristic is not two.
Conversely every such multiset gives the displayed solution.
Its partial-fraction residues distinguish the multisets. The
separant is $(z-X^2)R$, nonzero as a polynomial in $X$ at every
solution, so all solutions are regular.

For distinct prime-field labels, choose the roots to form a
multiplicative Sidon set including diagonal products. If $k$ roots
have already been chosen and $S$ is their $k(k+1)/2$ products, a
new root $x$ is excluded by at most
\[
 (k+2)|S|+k+1
\]
field elements: those with $xa\in S$ for an old root $a$, those
with $x^2\in S$, the old roots, and zero. Products among the new
pairs cannot collide otherwise. For $k\le N-1$ this bound is at
most $(N^3+N)/2$. A prime larger than this bound and $D+2$ therefore
suffices, and a prime of size $O(N^3)$ exists. This constructs all
the solutions and their pairwise distinct labels in the prime field.
\end{proof}

The solution count by itself therefore cannot prove linear
first-order MCA. The following agreement bound also shows why the
example is not a quadratic proximity-gap lower bound.

\begin{proposition}[Constant fixed-gap count for the isolated family]
\label{prop:isolated-family-mca}
Over a finite field satisfying Theorem~\ref{thm:quadratic-isolated-first-order},
fix a received line $f+zg$ on $n$ distinct coordinates and an agreement
threshold $D<A\le n$. Let $L$ be the number of distinct labels with
an $A$-agreeing candidate from the displayed family. Then
\[
 L\le \frac{(D+2)n}{A-D}.
\]
More strongly, let $r_0$ be the number of domain points that are
roots of $R$, put $m=n-r_0$ and $t=A-r_0$, and suppose $t^3>4m^2$.
Then
\[
 L\le \frac{2m}{t-(4m^2)^{1/3}}.
\]
In particular $L=O_\eta(1)$ for $A-D\ge\eta n$ and sufficiently
large $n$. Neither bound requires that all family labels be distinct.
\end{proposition}
\begin{proof}
Every candidate has degree $D$, hence at least $A-D$ of its
agreements have nonzero value. At a root $c$ of $R$, only the $N$
multisets containing $c$ have nonzero value. At a point $x$ outside
the roots, write $\alpha=f(x)/R(x)$ and $\beta=g(x)/R(x)$. For
fixed $a$, agreement is
\[
 \frac1{x-a}+\frac1{x-b}=\alpha+\beta ab.
\]
Clearing $x-b$ gives a nonzero polynomial of degree at most two
in $b$. If it vanished identically, its quadratic coefficient
would give $\beta=0$ since $a\ne0$, its linear coefficient would
give $\alpha=1/(x-a)$, and its constant coefficient would still
be nonzero. Thus there are at most $2N$ ordered matches. Diagonal
matches satisfy $2=(\alpha+\beta a^2)(x-a)$, a nonzero cubic
equation, so there are at most three of them. The number of
unordered matches is consequently at most
$\lfloor(2N+3)/2\rfloor=N+1$. Counting nonzero agreement incidences
proves the first bound, even for the number of nearby pairs.

For the stronger bound, choose one nearby candidate at each of
the $L$ distinct labels and put
\[
 W_{a,b}=P_{a,b}/R=\frac1{X-a}+\frac1{X-b}.
\]
No three of the labeled functions $(ab,W_{a,b})$ at distinct labels
are affine-collinear. To prove this, suppose their affine relation
has coefficients $c_1,c_2,c_3$, all nonzero, with
$\sum c_i=\sum c_i z_i=0$. A root belonging to only one support
cannot have its nonzero residue canceled. Each root in the union
therefore belongs to at least two supports, so the union has at
most three roots. If it has three, the supports must be the three
edges of a triangle. Their residue matrix has determinant $2$,
so they are linearly independent. If it has two roots $a,b$, the
three distinct multisets must be $\{a,a\},\{a,b\},\{b,b\}$.
Their only affine relation is
$W_{a,a}-2W_{a,b}+W_{b,b}=0$, but the same coefficients give
$a^2-2ab+b^2=(a-b)^2\ne0$ on the labels. One root cannot support
three distinct candidates. This proves the claim.

For any three distinct labels, take
$(c_1,c_2,c_3)=(z_2-z_3,z_3-z_1,z_1-z_2)$.
The nonzero rational function $\sum c_iW_i$ has a denominator of
degree at most six. Using the product of the three monic quadratic
denominators, its numerator has degree at most four: each summand
has leading term $2X^5$, which cancels because $\sum c_i=0$.
Thus three candidates can agree with the received line at at most
four common coordinates outside the roots of $R$.

Each selected candidate has at least $t$ agreements on these $m$
coordinates. If $h_x$ counts the agreeing candidates at $x$, then
$\sum h_x\ge tL$ and
\[
 \sum_x\binom{h_x}{3}\le4\binom L3\le\frac{4L^3}{6}.
\]
For every nonnegative integer $h$,
$6\binom h3\ge(h-2)_+^3$. Convexity therefore gives
\[
 m\left(\frac{tL}{m}-2\right)_+^3\le4L^3.
\]
Taking cube roots and rearranging proves the bound (the case
$tL\le2m$ also satisfies it). Finally $r_0\le D+1$, so
$t\ge\eta n-1$ and $m\le n$ at a fixed gap. The denominator is
positive and of order $n$ for sufficiently large $n$.
\end{proof}

\paragraph{Exact checks.}
The isolated-family checker exhausts $17{,}410$ polynomials,
covering $180{,}384$ polynomial--challenge pairs, and checks every
pair of local received values at each coordinate of four small-field
fixtures. It verifies the rational derivative identity and constructs
prime-field Sidon examples through $D=64$ with $2{,}145$ distinct
regular isolated labels. It also checks the nonzero degree-four
triple numerators and the integer convexity inequality. A characteristic
three negative control has an extra solution outside the theorem's
strict characteristic hypothesis. These finite checks supplement
the classification and incidence proofs above.
